\appendix

\section{Python Simulation Code}

Below is the Python code used to numerically simulate the computed force geodics of the golf ball along the golf course $h(q^1,q^2)$ for all examples and figures present in the paper involving golf ball trajectories.

\begin{lstlisting}
	
import numpy as np
import matplotlib.pyplot as plt
from mpl_toolkits.mplot3d import Axes3D
from matplotlib import cm
from scipy.integrate import solve_ivp
import sympy as sp


x_sym, y_sym = sp.symbols('x y')

# Pendulum case
r=3
eps=0.01
rho = sp.sqrt(x_sym**2 + y_sym**2)
u = rho / r
phi = (1 - u**2 + sp.sqrt((1 - u**2)**2 + eps**2))/2
h_expr =  - r * sp.sqrt(phi)

# Two planet gravity system
# h_expr = -1/((x_sym)**2+(y_sym-0.9)**2)-1/((x_sym)**2+y_sym**2)

#h_expr = 0.2*x_sym + 0.2*y_sym

# Compute derivatives symb.
hx_expr  = sp.diff(h_expr, x_sym)
hy_expr  = sp.diff(h_expr, y_sym)
hxx_expr = sp.diff(hx_expr, x_sym)
hyy_expr = sp.diff(hy_expr, y_sym)
hxy_expr = sp.diff(hx_expr, y_sym)

# Turn them into functions (through numpy which is fast)
h     = sp.lambdify((x_sym, y_sym), h_expr,  'numpy')
h_x   = sp.lambdify((x_sym, y_sym), hx_expr, 'numpy')
h_y   = sp.lambdify((x_sym, y_sym), hy_expr, 'numpy')
h_xx  = sp.lambdify((x_sym, y_sym), hxx_expr,'numpy')
h_yy  = sp.lambdify((x_sym, y_sym), hyy_expr,'numpy')
h_xy  = sp.lambdify((x_sym, y_sym), hxy_expr,'numpy')

# ---------------------------------
# Metric tensor and Christoffel symbols computation
# -----------------------------
def metric(x, y):
hx, hy = h_x(x,y), h_y(x,y)
G = np.array([[1+hx**2, hx*hy],
[hx*hy,   1+hy**2]])
return G

def christoffel(x, y):
# Christoffel symbols of first kind: Gamma_ijk = 1/2(partialj g_ik + partialk g_ij - partiali g_jk)
# Then raise index with inverse metric.
hx, hy = h_x(x,y), h_y(x,y)
hxx, hyy, hxy = h_xx(x,y), h_yy(x,y), h_xy(x,y)

G = metric(x,y)
Ginv = np.linalg.inv(G)

# Partial derivatives of metric
dGx = np.array([[2*hx*hxx, hxx*hy + hx*hxy],
[hxx*hy + hx*hxy, 2*hy*hxy]])
dGy = np.array([[2*hx*hxy, hxy*hy + hx*hyy],
[hxy*hy + hx*hyy, 2*hy*hyy]])

dG = [dGx, dGy]

# Compute Christoffel symbols Gamma^i_jk
Gamma = np.zeros((2,2,2))
for i in range(2):
for j in range(2):
for k in range(2):
term = 0.0
for l in range(2):
term += Ginv[i,l]*(dG[j][l,k] + dG[k][l,j] - dG[l][j,k])
Gamma[i,j,k] = 0.5*term
return Gamma

# -----------------------------
# Equations of motion
# -----------------------------
g = 9.81
m = 1.0

def dynamics(t, state):
x, y, vx, vy = state
qdot = np.array([vx, vy])

G = metric(x,y)
Ginv = np.linalg.inv(G)
Gamma = christoffel(x,y)

# Compute RHS force term: -G^{ij} partialj V/m, with V = m g h(x,y)
gradV = np.array([m*g*h_x(x,y), m*g*h_y(x,y)])
force = -Ginv @ (gradV/m)

# Acceleration: -Gamma^i_jk v^j v^k + force^i
accel = np.zeros(2)
for i in range(2):
for j in range(2):
for k in range(2):
accel[i] -= Gamma[i,j,k]*qdot[j]*qdot[k]
accel[i] += force[i]

return [vx, vy, accel[0], accel[1]]

# -----------------------------
# Initial conditions
# -----------------------------
x0, y0 = 1, 2
vx0, vy0 = -3, 5
state0 = [x0, y0, vx0, vy0]

# Time span

t_span = (0, 100) # n seconds
t_eval = np.linspace(t_span[0], t_span[1], 1000)

sol = solve_ivp(dynamics, t_span, state0, t_eval=t_eval, rtol=1e-8, atol=1e-8)

# Extract trajectory
X, Y = sol.y[0], sol.y[1]
Z = h(X,Y)
import plotly.graph_objects as go
import numpy as np

# Create surface
gx = np.linspace(-4, 4, 500)
gy = np.linspace(-4, 4, 500)
GX, GY = np.meshgrid(gx, gy)
GZ = h(GX, GY)

# Surface
surface = go.Surface(x=GX, y=GY, z=GZ, colorscale="Greens", opacity=0.9)

# Trajectory
traj = go.Scatter3d(x=X, y=Y, z=Z, mode="lines", line=dict(color="red", width=5), name="Trajectory")

# Start/End points
start = go.Scatter3d(x=[X[0]], y=[Y[0]], z=[Z[0]], mode="markers", marker=dict(size=6, color="blue"), name="Start")
end   = go.Scatter3d(x=[X[-1]], y=[Y[-1]], z=[Z[-1]], mode="markers", marker=dict(size=6, color="black"), name="End")

fig = go.Figure(data=[surface, traj, start, end])
fig.update_layout(scene=dict(
xaxis_title="x",
yaxis_title="y",
zaxis_title="z = h(x,y)",
#zaxis=dict(range=[-1, 1])
))
fig.show()
\end{lstlisting}

\phantomsection
\addcontentsline{toc}{section}{References}
